
\section{奖励与荣誉}

\datedline{\textbf{Data-电商2024 H2 \& Q4 Spot Bonus 超出预期结果奖}}{2024}
字节跳动，Data-电商团队，供应链与物流算法部门100多人团队中，仅4人获奖
\vspace{5pt}

\datedline{\textbf{入选国家级青年人才项目（海外）}}{2024}
依托清华大学土木工程系申请
\vspace{5pt}


\datedline{\textbf{Global Urban Rail Transit Best Dissertation Award \& Zhongheng Shi Honorary Prize}}{2023}
第一名，每年奖励给一个在轨道交通领域最优秀的博士毕业论文，被 {\href{https://link.springer.com/journal/40864/updates/27693480#:~:text=The%20Journal%20of%20Urban%20Rail,field%20of%20Urban%20Rail%20Transit.}{\textcolor{blue}{{ The Journal of Urban Rail Transit}}}} 报道
\vspace{5pt}


\datedline{\textbf{Dan \& Eva Roos Thesis Prize, MIT}}{2023}
每年奖励给一个MIT在交通、物流领域最优秀的博士毕业论文（申请者来自MIT计算机系、城市规划系、经济系、运筹学中心等），在MIT Mobility Forum进行展示
\vspace{5pt}

\datedline{\textbf{COTA最佳博士论文奖}}{2022}
每年授予一至三名交通领域最优秀的博士毕业生，在TRB会议进行展示
\vspace{5pt}

\datedline{\textbf{Amazon最后一公里路径规划挑战赛，全球第二名}}{2021}
\$50,000现金奖励，被 {\href{https://www.amazon.science/academic-engagements/winning-last-mile-challenge-team-addresses-problem-of-combining-mathematical-routes-with-driver-knowledge}{\textcolor{blue}{{Amazon Science}}}} 和   {\href{https://news.mit.edu/2021/last-mile-routing-research-challenge-three-winning-teams-0824}{\textcolor{blue}{{MIT News}}}} 报道，唯一学生获奖队伍，第一名为三位滑铁卢大学数学系终身教授
\vspace{5pt}

\datedline{\textbf{UPS PhD Fellowship, MIT}}{2021}
\$90,000奖学金，每年奖励给一名MIT在交通、物流领域最优秀的博士生（申请者来自MIT计算机系、城市规划系、经济系、运筹学中心等）{\href{https://ctl.mit.edu/news/ups-fellowships-maintain-tradition-graduate-student-support}{\textcolor{blue}{{（报道链接）}}}} 
\vspace{5pt}

\datedline{\textbf{清华大学优秀毕业论文}}{2018}
前3\%的本科生优秀毕业论文
\vspace{5pt}

\datedline{\textbf{清华大学优秀毕业生、北京市优秀毕业生}}{2018}
前1\%的本科优秀毕业生
\vspace{5pt}

\datedline{\textbf{清华大学本科生特等奖学金}}{2017}
清华大学本科生最高荣誉，每年从所有本科生中评选10人，评选过程得到人民日报报道并上微博热搜，评选最后一轮面向全校学生公开答辩，在学校官网直播{\href{http://www.tuef.tsinghua.edu.cn/info/qhyx/3216}{\textcolor{blue}{{（报道链接）}}}}
\vspace{5pt}

\datedline{\textbf{清华大学蔡雄奖学金}}{2017}
奖励给科创拔尖的优秀本科生，每年从所有本科生中评选10人，作为10人代表进行感谢发言
\vspace{5pt}

\datedline{\textbf{清华大学本科生学术推进计划}}{2017}
支持优秀本科生自我立项进行独立科研，获得最高级评定，作为优秀代表与薛其坤副校长交流，并在学术推进计划年会进行项目展示
\vspace{5pt}

\datedline{\textbf{清华大学结构设计大赛特等奖}}{2016}
担任队长，包揽所有单项奖，清华大学三星级赛事（全校仅有5个），土木工程系最重要的比赛
\vspace{5pt}

\datedline{\textbf{清华大学唐立新奖学金}}{2016 -- 2022}
奖励给学业、科创、社工等综合能力突出的学生，清华每年约30人，一直奖励到博士毕业，是清华当年唯二在大二年级获此奖项的本科生
\vspace{5pt}

\datedline{\textbf{清华大学闯世界计划}}{2016，2017}
支持优秀本科生进行海外研修，连续两年（2016、2017）获得最高级（A级）评定，是2016年唯一获得A级评定的大二年级本科生
\vspace{5pt}

\datedline{\textbf{清华大学星火计划}}{2016 -- 2018}
清华大学人才培养计划，致力于支持优秀本科生进行独立科研活动，每年从全校选拔约40人
\vspace{5pt}

\datedline{\textbf{国家奖学金}}{2015}
中华人民共和国教育部给本科生的最高奖学金，约前0.1\%
\vspace{5pt}


\datedline{\textbf{清华大学唐仲英德育奖学金}}{2015 -- 2018}
奖励给学业优秀和有志于社会公益的优秀本科生，每年约40人，一直奖励到本科毕业
\vspace{5pt}

\datedline{\textbf{清华大学综合优秀奖、学业优秀奖、科创竞赛奖、社会工作奖、优秀共青团员、优秀学生干部、三星级志愿者}}{2015 -- 2018}
在读期间获得清华大学多项荣誉
\vspace{5pt}

\datedline{\textbf{2014年北京市高考第11名、北京市三好学生}}{2014}
高考分数712/750（含10分北京市三好学生加分，北京市前0.02\%），被清华大学录取
